\section{Benchmark audit}
\label{sec:findings}

The findings in this section do not depend on our method. They concern the dataset, its
standard split and the conventions under which group-level results are scored, and they affect
any result published on the AMLworld family.

\subsection{The test tail}
\label{sec:findings-tail}

The day split places a pathology in the test tail (Figure~\ref{fig:tail}). Days 10--17 hold
655 of the test period's 1{,}611 positives inside 1{,}108 edges --- a 59\,\% base rate against
0.11\,\% on days 8--9 --- so 40.7\,\% of the test positives sit in 0.13\,\% of the test edges,
and a rule that reads the clock (``after day 10 $\Rightarrow$ illicit'') scores F1 0.48 there
on its own.
% source: docs/FINDINGS.md (C1); docs/TEST_TAIL.md
This is a property of the generator, not of one release: regenerating the full preprocessing
path on LI-Small reproduces the collapse (Table~\ref{tab:tail}). What reproduces is the
concentration; the share of test signal in the tail differs by release.
% source: docs/paper/evaluation_protocol.md §4

\begin{figure}[t]
\centering
\includegraphics[width=0.62\textwidth]{fig_tail_base_rate.png}
\caption{Laundering base rate by test day in two AMLworld releases. After day 9 the
background traffic collapses and the base rate jumps by more than two orders of magnitude.}
% source: docs/paper/evaluation_protocol.md §4 (per-day base-rate table); figure built by paper_draft/make_figs.py
\label{fig:tail}
\end{figure}

% Source: docs/paper/evaluation_protocol.md §4 ("This is a property of the generator, not of HI-Small")
\begin{table}[t]
\centering
\caption{The day-split test tail in two AMLworld releases. The concentration reproduces
(${\le}\,0.13\,\%$ of test edges at a ${\sim}60\,\%$ base rate); the share of test signal in
the tail differs by release.}
\label{tab:tail}
\begin{tabular}{lrr}
\toprule
 & HI-Small & LI-Small \\
\midrule
Tail edges (days $\ge 10$) & 1{,}108 (0.128\,\% of test) & 223 (0.019\,\% of test) \\ % source: docs/paper/evaluation_protocol.md §4 (table)
Test positives in the tail & 655 (40.7\,\%) & 134 (16.7\,\%) \\ % source: docs/paper/evaluation_protocol.md §4 (table)
Tail base rate & 59.1\,\% & 60.1\,\% \\ % source: docs/paper/evaluation_protocol.md §4 (table)
Lift over the split's test base rate & 318$\times$ & 880$\times$ \\ % source: docs/paper/evaluation_protocol.md §4 (table)
\bottomrule
\end{tabular}
\end{table}


Two consequences follow. A single-number test F1 on HI-Small is not comparable across papers
unless the window is stated: our detector's edge F1 is $0.495 \pm 0.027$ on days 8--9 and
0.596 on the whole split.
% source: results/tables/r1_edge_metrics.csv (primary_f1_minority_mean/std); r6_detector_swap.csv (edge_f1, gine_causal)
And the tail reaches ring-level results through the seed budget. A global 0.1\,\% score
quantile over the test split places only $463.0$ of its 864 seeds in the two clean days, so
any method that selects by a global quantile over this split operates in the clean window at
roughly half its nominal budget (Section~\ref{sec:results-swap} measures what recovering that
budget does).
% source: results/tables/v34_window_budget.json; docs/FINDINGS.md (dated note, window-matched seed budget)

\subsection{The day-0 warm-up}

Day 0 holds 34.3\,\% of HI-Small's training edges (34.4\,\% of LI-Small's) at roughly ten
times day 1's per-minute density. Trailing motif features count structures inside a time
window, so that density distorts them in both directions: cycle-shaped features are nonzero on
12.6\,\% of day-0 edges against about 2\,\% on every later day, while fan-in and fan-out
features are depressed (63\,\% against 89\,\%) because there is no history to look back on.
Both effects reproduce in LI-Small to within 0.3 percentage points.
% source: docs/FINDINGS.md (C8)
Day 0 is the simulation warming up, and anything trained across it is partly fitting that
warm-up.

How much this matters is release-specific. On HI-Medium, where day 0 is 48.8\,\% of training
edges at 0.27 times the later base rate, a detector trained across it reaches validation AP
0.3073; excluding day-0 \emph{labels} while keeping the day as message-passing history gives
0.4495 on half the supervised edges.
% source: docs/FINDINGS.md (C10)
The same exclusion does not help on HI-Small: with a gate fixed before training, seed-1
validation AP is 0.5455 against the reference detector's 0.5726, and the remaining seeds were
not run.
% source: docs/deep-analysis.md §9 — one seed; not yet recorded in docs/FINDINGS.md

\subsection{A global quantile can hand the seed budget to one day}

Every release ends its ordinary traffic with an elevated-volume day (1.36--1.57 times the
weekday norm). On HI-Medium, a single global 0.1\,\% quantile over the test split gives that
day 87.5\,\% of the seeds although it holds 16\,\% of the window's edges, because the
detector's scores are systematically higher there. On HI-Small the corresponding day takes
22.0--36.3\,\% of seeds across detector seeds against 75.8\,\% of the split's edges, so volume
alone does not explain the capture.
% source: docs/FINDINGS.md (C11); results/tables/v33_seed_by_day.csv
Budget allocation across days is a design choice and should be stated, not left to the score
distribution.

\subsection{Saturated scores break strict quantile cuts}
\label{sec:findings-saturation}

On one HI-Medium detector seed, 780{,}093 of 18.9\,M test scores are exactly 1.0. The 99.9\,\%
quantile is then 1.0 itself, and a strict ``score $>$ q'' rule selects no seeds at all.
Saturation grows with distance from training: 124 saturated edges on test day 6, 25{,}495 on
day 9 and 221{,}909 on day 15. Results from such a seed are an artefact of the cut, not a
measurement, and we exclude that seed from every HI-Medium number. The general point is that
a quantile rule is only defined once its tie handling is stated.
% source: docs/deep-analysis.md §7 (diagnosis 1: seed 3, 780,093 of 18.9 M scores at 1.0; 124 / 25,495 / 221,909 on days 6 / 9 / 15)

\subsection{The scope convention is worth five times the method}
\label{sec:findings-scope}

Figure~\ref{fig:scope} re-scores the same runs under four rules for what a truth ring's member
set is. Nothing is re-run and no parameter is re-chosen.

\begin{figure}[t]
\centering
\includegraphics[width=0.78\textwidth]{fig_scope_conventions.png}
\caption{Ring-F1 of the same runs under four scope conventions (Multi-GINe detector, primary
window). The convention moves every method by far more than the methods differ from each
other.}
% source: results/tables/r14_scope_sensitivity.csv and v36_component_scope.csv; figure built by paper_draft/make_figs.py
\label{fig:scope}
\end{figure}

Counting all of a ring's accounts as members, not only those visible in the window, lowers
ring-F1 by 0.347--0.383 for every method and both GNN detectors (Table~\ref{tab:scope},
appendix). The in-scope ring count is unchanged at 174; only member sets grow, from a median
of 4 accounts to 11. For scale, the largest method effect in this paper is $+0.068$.
% source: results/tables/r14_scope_sensitivity.csv; docs/FINDINGS.md (C9)

Making targets smaller instead --- scoring each in-window connected piece of a ring as its own
target --- gives 254 targets; keeping only pieces of three or more accounts gives 108 targets
in 97 rings (Table~\ref{tab:component-scope}, appendix). The ordering
seeds $<$ rule $<$ policy $<$ policy-plus-partition is preserved under three of the four
conventions for both detectors. Under the fourth it changes in one place: the policy alone
overtakes policy-plus-partition (0.766 vs.\ 0.740). The swap was not tested for significance,
but the sign of the partition step flips, and it flips again on HI-Medium
(Section~\ref{sec:findings-hm}).
% source: results/tables/v36_component_scope.csv; docs/deep-analysis.md §4
Adding the community-detection baselines of Section~\ref{sec:results-baselines} shows what
re-ranking by convention looks like beyond that swap. Conventions re-rank methods that are
close, not methods that are clearly apart. The only other swaps between neighbours in the
ranking are seeds-plus-components against Louvain at resolution 0.5 (0.4436 vs.\ 0.4451 when
each in-window piece is a target; 0.6731 vs.\ 0.6738 with pieces of three or more accounts)
and the two weakest methods, Louvain at resolution 2 against CSGM, when all of a ring's
accounts count.
% source: results/tables/v40_external_baselines.json (bar_A adjacent_swaps); v40_external_baselines.csv (b0 / louvain_r0.5 component and component_ge3 ring_f1; louvain_r2 / csgm full_membership)

Neither extreme is right. Scoring against in-window accounts holds a method responsible only
for evidence it could have seen, and makes patterns severed by the split structurally
unmatchable. Scoring against all accounts keeps rings intact and penalises a perfect in-window
method for members it could not observe. A ring-level number is not interpretable without
knowing which rule produced it, and none of the group-level work we reviewed in
Section~\ref{sec:related} states one.
% source: docs/paper/evaluation_protocol.md §9

\subsection{A second release with ring-level evaluation: HI-Medium}
\label{sec:findings-hm}

HI-Medium has 2{,}187 in-scope rings in its primary window (test days 6--14), against
HI-Small's 174. We re-scored its existing ring outputs under the same protocol, with no new
training.
% source: results/tables/v41_hm_audit.json (observability.n_in_scope_rings; window); results/notes/p4_hm_audit.md
\emph{The direction of the scope effect reproduces; its size does not.} Counting all of a
ring's accounts lowers ring-F1 for every method under both seed-allocation rules, by 0.038 to
0.059 under the global rule and 0.038 to 0.073 across all eight cells
(Table~\ref{tab:hm-audit}). That is smaller than the difference between the full pipeline and
seeds-plus-components on the same release (0.092), where on HI-Small the scope gap was five
times the method effect. One hypothesis, which we have not tested, is window length: HI-Medium's
primary window spans nine days against HI-Small's two, so more of each ring is visible and
restricting members to the window removes less.
% source: results/tables/v41_hm_audit_scope.csv (gap_f1_fm_minus_tv); v41_hm_audit.json (verdicts.scope.global: largest_method_effect_full_minus_b0 0.0924)
\emph{Observability reproduces.} 875 of the 2{,}187 rings (40.0\,\%) have no connected in-window
piece of three or more accounts, against 44.3\,\% on HI-Small.
% source: results/tables/v41_hm_audit.json (observability: rings_no_in_window_piece_ge3 875, n_in_scope_rings 2187)
\emph{The threshold dependence of the partition gain reproduces.} The gain in ring recall from
partitioning falls from 0.019 at Jaccard 0.3 to 0.015 at 0.5 and 0.009 at 0.7 (global rule),
and under the scope that keeps only in-window pieces of three or more accounts the policy alone
again edges past policy-plus-partition (0.2434 vs.\ 0.2430).
% source: results/tables/v41_hm_audit_threshold.csv (global gain_ring_recall 0.0194 / 0.0149 / 0.0089); v41_hm_audit_component.csv (global proposed / full ring_f1_component_ge3)

% Source: results/tables/v41_hm_audit_scope_v2.csv (tv_ring_f1, fm_ring_f1, gap_f1_fm_minus_tv), point estimates only;
% points identical to v41_hm_audit_scope.csv (only the intervals were re-scored).
\begin{table}[t]
\centering
\caption{The scope gap on HI-Medium (2{,}187 in-scope rings, test days 6--14, two detector
seeds, Jaccard $\ge 0.5$). Ring-F1 when a ring's members are its in-window accounts, when they
are all its accounts, and the difference. Point estimates over the two seeds.$^{\ddagger}$}
\label{tab:hm-audit}
\small
\begin{tabular}{llccc}
\toprule
Seed allocation & Method & In-window members & All members & Gap \\
\midrule
global  & seeds + components       & 0.090 & 0.052 & $-0.038$ \\ % source: v41_hm_audit_scope_v2.csv (global,b0)
global  & rule expansion           & 0.097 & 0.054 & $-0.043$ \\ % source: v41_hm_audit_scope_v2.csv (global,b1b)
global  & learned expansion        & 0.180 & 0.121 & $-0.059$ \\ % source: v41_hm_audit_scope_v2.csv (global,proposed)
global  & expansion + partition    & 0.182 & 0.126 & $-0.057$ \\ % source: v41_hm_audit_scope_v2.csv (global,full)
\midrule
per day & seeds + components       & 0.148 & 0.087 & $-0.061$ \\ % source: v41_hm_audit_scope_v2.csv (per_day,b0)
per day & rule expansion           & 0.141 & 0.082 & $-0.059$ \\ % source: v41_hm_audit_scope_v2.csv (per_day,b1b)
per day & learned expansion        & 0.233 & 0.160 & $-0.073$ \\ % source: v41_hm_audit_scope_v2.csv (per_day,proposed)
per day & expansion + partition    & 0.234 & 0.165 & $-0.069$ \\ % source: v41_hm_audit_scope_v2.csv (per_day,full)
\bottomrule
\multicolumn{5}{l}{\footnotesize $^{\ddagger}$ HI-Medium: detector seeds 1--2 only (seed 3's scores saturate); the detector did not}\\
\multicolumn{5}{l}{\footnotesize meet its pre-registered adequacy bar, so these rows are an audit replication, not a method comparison.}\\
\end{tabular}
\end{table}


\subsection{A release without enough rings: LI-Small}

LI-Small was regenerated end to end (6{,}924{,}049 transactions, 705{,}907 accounts). The
preprocessing, splitting, scoping and windowing path runs unchanged, and the tail and day-0
artefacts reproduce. Its ring-level evaluation is under-powered --- 38 matchable rings in its
primary window against HI-Small's 114, below a stopping threshold fixed in advance --- so we
report no method numbers on it.
% source: docs/paper/evaluation_protocol.md §4 (LI-Small subsection)
